% FILE: CSE-23_CT-1.tex
\documentclass[11pt, a4paper]{article}
\usepackage[a4paper, top=2.5cm, bottom=2.5cm, left=2cm, right=2cm]{geometry}
\usepackage{fontspec}
\usepackage[english]{babel}
\usepackage{amsmath}
\usepackage{amssymb}
\usepackage{enumitem}

\babelfont{rm}{Noto Sans}

\begin{document}

%%%%%%%%%%%%%%%%%%%%%%%%%%%%%%%%%%%%%%%%%%%%%%%%%%  
% QUESTION_START  
%%%%%%%%%%%%%%%%%%%%%%%%%%%%%%%%%%%%%%%%%%%%%%%%%%
% id=cse23-ct1-seca-q1  
% discipline=CSE  
% year=23  
% exam_type=CT  
% section=A  
% ct_number=1  
% question_number=1  
% topic=Functions  
% subtopic=Curve Sketching  
% difficulty=Medium  
% length=Medium  
% question_type=New  
% frequency=1  
% appearances=  
% OCR_STATUS=VERIFIED

\section*{CT-1 (Sec A) - Q1}

Question:
Draw the graph of the following function and find the domain and range of it:
\[ 
f(x) = \begin{cases} 
0, & x \le -1 \\ 
\sqrt{1-x^2}, & -1 < x < 1 \\ 
x, & x \ge 1 
\end{cases} 
\]

Solution:

\begin{description}
    \item[Step 1: Determine the Domain] \ \\
    The function piecewise branches cover the entire real line:
    \begin{itemize}
        \item $x \le -1$
        \item $-1 < x < 1$
        \item $x \ge 1$
    \end{itemize}
    Combining these intervals, the domain of the function is:
    \[ \text{Domain} = (-\infty, \infty) \]

    \item[Step 2: Determine the Range] \ \\
    We analyze the range of each piecewise branch:
    \begin{enumerate}
        \item For $x \le -1$:
        \[ f(x) = 0 \implies f(x) \in \{0\} \]
        \item For $-1 < x < 1$:
        \[ f(x) = \sqrt{1-x^2} \]
        Since $0 \le x^2 < 1$, we have $0 < 1 - x^2 \le 1$, which gives:
        \[ 0 < \sqrt{1-x^2} \le 1 \implies f(x) \in (0, 1] \]
        \item For $x \ge 1$:
        \[ f(x) = x \implies f(x) \in [1, \infty) \]
    \end{enumerate}
    Combining the outputs from all three branches:
    \[ \text{Range} = \{0\} \cup (0, 1] \cup [1, \infty) = [0, \infty) \]

    \item[Step 3: Graph Description] \ 
    \begin{itemize}
        \item For $x \in (-\infty, -1]$, the graph is a flat horizontal line lying directly on the x-axis ($y = 0$).
        \item For $x \in (-1, 1)$, the graph is the upper half of a unit circle centered at the origin, starting with an open circle at $(-1, 0)$, reaching its peak at $(0, 1)$, and ending with an open circle at $(1, 0)$.
        \item For $x \in [1, \infty)$, the graph is a straight line $y = x$ starting at the solid point $(1, 1)$ and extending upwards to infinity at a $45^\circ$ angle.
    \end{itemize}
\end{description}

Final Answer:
\[ 
\boxed{\text{Domain} = (-\infty, \infty), \quad \text{Range} = [0, \infty)} 
\]

%%%%%%%%%%%%%%%%%%%%%%%%%%%%%%%%%%%%%%%%%%%%%%%%%%  
% QUESTION_END  
%%%%%%%%%%%%%%%%%%%%%%%%%%%%%%%%%%%%%%%%%%%%%%%%%%

%%%%%%%%%%%%%%%%%%%%%%%%%%%%%%%%%%%%%%%%%%%%%%%%%%  
% QUESTION_START  
%%%%%%%%%%%%%%%%%%%%%%%%%%%%%%%%%%%%%%%%%%%%%%%%%%
% id=cse23-ct1-seca-q2  
% discipline=CSE  
% year=23  
% exam_type=CT  
% section=A  
% ct_number=1  
% question_number=2  
% topic=Continuity  
% subtopic=General  
% difficulty=Medium  
% length=Medium  
% question_type=New  
% frequency=1  
% appearances=  
% OCR_STATUS=VERIFIED

\section*{CT-1 (Sec A) - Q2}

Question:
Discuss the continuity of the following function at $x = 0, 3/2$:
\[ 
f(x) = \begin{cases} 
3+2x, & \text{for } -3/2 \le x \le 0 \\ 
3-2x, & \text{for } 0 \le x < 3/2 \\ 
-3-2x, & \text{for } x \ge 3/2 
\end{cases} 
\]

Solution:

\begin{description}
    \item[Step 1: Test continuity at $x = 0$] \ \\
    A function is continuous at $x = 0$ if and only if $\lim_{x \to 0^-} f(x) = \lim_{x \to 0^+} f(x) = f(0)$.
    \begin{itemize}
        \item \textbf{Function value:} $f(0) = 3 + 2(0) = 3$ (using the first branch).
        \item \textbf{Left-Hand Limit (LHL):}
        \[ \lim_{x \to 0^-} f(x) = \lim_{x \to 0^-} (3 + 2x) = 3 \]
        \item \textbf{Right-Hand Limit (RHL):}
        \[ \lim_{x \to 0^+} f(x) = \lim_{x \to 0^+} (3 - 2x) = 3 \]
    \end{itemize}
    Since $\text{LHL} = \text{RHL} = f(0) = 3$, the function $f(x)$ is \textbf{continuous} at $x = 0$.

    \item[Step 2: Test continuity at $x = 2$] \ \\
    We check if $\lim_{x \to 2^-} f(x) = \lim_{x \to 2^+} f(x) = f(2)$.
    \begin{itemize}
        \item \textbf{Function value:} $f(3/2) = -3 - 2(3/2) = -3 - 3 = -6$ (using the third branch).
        \item \textbf{Left-Hand Limit (LHL):}
        \[ \lim_{x \to 3/2^-} f(x) = \lim_{x \to 3/2^-} (3 - 2x) = 3 - 2\left(\frac{3}{2}\right) = 0 \]
        \item \textbf{Right-Hand Limit (RHL):}
        \[ \lim_{x \to 3/2^+} f(x) = \lim_{x \to 3/2^+} (-3 - 2x) = -3 - 2\left(\frac{3}{2}\right) = -6 \]
    \end{itemize}
    Since $\text{LHL} \neq \text{RHL}$ ($0 \neq -6$), the limit at $x = 3/2$ does not exist. Hence, $f(x)$ has a jump discontinuity and is \textbf{discontinuous} at $x = 3/2$.
\end{description}

Final Answer:
\[ 
\boxed{f(x) \text{ is continuous at } x = 0 \text{ and discontinuous at } x = 3/2.} 
\]

%%%%%%%%%%%%%%%%%%%%%%%%%%%%%%%%%%%%%%%%%%%%%%%%%%  
% QUESTION_END  
%%%%%%%%%%%%%%%%%%%%%%%%%%%%%%%%%%%%%%%%%%%%%%%%%%

%%%%%%%%%%%%%%%%%%%%%%%%%%%%%%%%%%%%%%%%%%%%%%%%%%  
% QUESTION_START  
%%%%%%%%%%%%%%%%%%%%%%%%%%%%%%%%%%%%%%%%%%%%%%%%%%
% id=cse23-ct1-seca-q3  
% discipline=CSE  
% year=23  
% exam_type=CT  
% section=A  
% ct_number=1  
% question_number=3  
% topic=Differentiation  
% subtopic=Implicit Differentiation  
% difficulty=Medium  
% length=Long  
% question_type=New  
% frequency=1  
% appearances=  
% OCR_STATUS=VERIFIED

\section*{CT-1 (Sec A) - Q3}

Question:
Define explicit function and implicit function. Find $dy/dx$ of:
\begin{enumerate}[label=(\roman*)]
    \item $x^y + y^x = a^b$
    \item $y = x^{\tan^{-1}x} + (\cos x)^{\ln x}$
    \item $x = e^t \sin t, \ y = e^t \cos t$
\end{enumerate}

Solution:

\begin{description}
    \item[Definitions:] \ 
    \begin{itemize}
        \item \textbf{Explicit Function:} A function in which the dependent variable $y$ is expressed directly and solely in terms of the independent variable $x$. It is written in the form $y = f(x)$. \\
        \textit{Example:} $y = x^2 + 5x + 6$.
        \item \textbf{Implicit Function:} A function in which the dependent variable $y$ is not expressed directly in terms of $x$, but is instead defined via an equation involving both variables. It is written in the form $F(x, y) = 0$. \\
        \textit{Example:} $x^2 + y^2 - 25 = 0$.
    \end{itemize}
\end{description}

\subsection*{(i) Find $\frac{dy}{dx}$ of $x^y + y^x = a^b$}
\begin{description}
    \item[Step 1: Use substitution] \ \\
    Let $u = x^y$ and $v = y^x$. Thus:
    \[ u + v = a^b \implies \frac{du}{dx} + \frac{dv}{dx} = 0 \]

    \item[Step 2: Differentiate $u = x^y$] \ \\
    Taking natural logs:
    \[ \ln u = y \ln x \]
    Differentiating implicitly with respect to $x$:
    \[ \frac{1}{u}\frac{du}{dx} = y_1 \ln x + \frac{y}{x} \implies \frac{du}{dx} = x^y \left( y_1 \ln x + \frac{y}{x} \right) \]

    \item[Step 3: Differentiate $v = y^x$] \ \\
    Taking natural logs:
    \[ \ln v = x \ln y \]
    Differentiating implicitly with respect to $x$:
    \[ \frac{1}{v}\frac{dv}{dx} = \ln y + \frac{x}{y} y_1 \implies \frac{dv}{dx} = y^x \left( \ln y + \frac{x}{y} y_1 \right) \]

    \item[Step 4: Combine and solve for $y_1$] \ \\
    \[ x^y \left( y_1 \ln x + \frac{y}{x} \right) + y^x \left( \ln y + \frac{x}{y} y_1 \right) = 0 \]
    \[ y_1 \left( x^y \ln x + x y^{x-1} \right) + y x^{y-1} + y^x \ln y = 0 \]
    \[ y_1 = -\frac{y x^{y-1} + y^x \ln y}{x^y \ln x + x y^{x-1}} \]
\end{description}

\subsection*{(ii) Find $\frac{dy}{dx}$ of $y = x^{\tan^{-1}x} + (\cos x)^{\ln x}$}
\begin{description}
    \item[Step 1: Differentiate each term independently] \ \\
    Let $y = u + v$, where $u = x^{\tan^{-1}x}$ and $v = (\cos x)^{\ln x}$.
    \begin{itemize}
        \item \textbf{For $u = x^{\tan^{-1}x}$:}
        \[ \ln u = \tan^{-1}x \ln x \implies \frac{1}{u}\frac{du}{dx} = \frac{\ln x}{1+x^2} + \frac{\tan^{-1}x}{x} \]
        \[ \frac{du}{dx} = x^{\tan^{-1}x} \left( \frac{\ln x}{1+x^2} + \frac{\tan^{-1}x}{x} \right) \]
        \item \textbf{For $v = (\cos x)^{\ln x}$:}
        \[ \ln v = \ln x \ln(\cos x) \implies \frac{1}{v}\frac{dv}{dx} = \frac{\ln(\cos x)}{x} + \ln x \left( -\frac{\sin x}{\cos x} \right) \]
        \[ \frac{dv}{dx} = (\cos x)^{\ln x} \left( \frac{\ln(\cos x)}{x} - \ln x \tan x \right) \]
    \end{itemize}
    Combining these derivatives:
    \[ \frac{dy}{dx} = x^{\tan^{-1}x} \left( \frac{\ln x}{1+x^2} + \frac{\tan^{-1}x}{x} \right) + (\cos x)^{\ln x} \left( \frac{\ln(\cos x)}{x} - \ln x \tan x \right) \]
\end{description}

\subsection*{(iii) Find $\frac{dy}{dx}$ of $x = e^t \sin t, \ y = e^t \cos t$}
\begin{description}
    \item[Step 1: Perform parametric differentiation] \ \\
    Differentiating $x$ with respect to $t$:
    \[ \frac{dx}{dt} = e^t \sin t + e^t \cos t = e^t(\sin t + \cos t) \]
    Differentiating $y$ with respect to $t$:
    \[ \frac{dy}{dt} = e^t \cos t - e^t \sin t = e^t(\cos t - \sin t) \]
    Finding $\frac{dy}{dx}$ using the chain rule:
    \[ \frac{dy}{dx} = \frac{dy/dt}{dx/dt} = \frac{e^t(\cos t - \sin t)}{e^t(\sin t + \cos t)} = \frac{\cos t - \sin t}{\sin t + \cos t} \]
\end{description}

Final Answer:
\[ 
\boxed{
\begin{aligned}
&\text{(i) } \frac{dy}{dx} = -\frac{y x^{y-1} + y^x \ln y}{x^y \ln x + x y^{x-1}} \\
&\text{(ii) } \frac{dy}{dx} = x^{\tan^{-1}x} \left( \frac{\ln x}{1+x^2} + \frac{\tan^{-1}x}{x} \right) + (\cos x)^{\ln x} \left( \frac{\ln(\cos x)}{x} - \ln x \tan x \right) \\
&\text{(iii) } \frac{dy}{dx} = \frac{\cos t - \sin t}{\sin t + \cos t}
\end{aligned}
} 
\]

%%%%%%%%%%%%%%%%%%%%%%%%%%%%%%%%%%%%%%%%%%%%%%%%%%  
% QUESTION_END  
%%%%%%%%%%%%%%%%%%%%%%%%%%%%%%%%%%%%%%%%%%%%%%%%%%

%%%%%%%%%%%%%%%%%%%%%%%%%%%%%%%%%%%%%%%%%%%%%%%%%%  
% QUESTION_START  
%%%%%%%%%%%%%%%%%%%%%%%%%%%%%%%%%%%%%%%%%%%%%%%%%%
% id=cse23-ct1-seca-q4  
% discipline=CSE  
% year=23  
% exam_type=CT  
% section=A  
% ct_number=1  
% question_number=4  
% topic=Differentiation  
% subtopic=Leibnitz Rule  
% difficulty=Hard  
% length=Long  
% question_type=New  
% frequency=1  
% appearances=  
% OCR_STATUS=VERIFIED

\section*{CT-1 (Sec A) - Q4}

Question:
State and prove Leibnitz's theorem. If $y = (\sin^{-1} x)^2$ then show that 
\[ (1-x^2)y_{n+2} - (2n+1)xy_{n+1} - n^2 y_n = 0 \]
Also find $(y_n)_0$.

Solution:

\begin{description}
    \item[Step 1: State Leibnitz's Theorem] \ \\
    If $u$ and $v$ are two functions of $x$ that can be differentiated $n$ times, then the $n$-th order derivative of their product $uv$ is given by:
    \[ (uv)_n = u_n v + \binom{n}{1} u_{n-1} v_1 + \binom{n}{2} u_{n-2} v_2 + \dots + \binom{n}{r} u_{n-r} v_r + \dots + u v_n \]

    \item[Step 2: Prove Leibnitz's Theorem] \ \\
    We prove the theorem using Mathematical Induction.
    \begin{itemize}
        \item \textbf{Base Case ($n=1$):} By the standard product rule:
        \[ (uv)_1 = u_1 v + u v_1 \]
        Since $\binom{1}{0}=1$ and $\binom{1}{1}=1$, the theorem holds for $n=1$.
        \item \textbf{Inductive Hypothesis:} Assume the theorem is true for $n=k$:
        \[ (uv)_k = u_k v + \binom{k}{1} u_{k-1} v_1 + \binom{k}{2} u_{k-2} v_2 + \dots + \binom{k}{r} u_{k-r} v_r + \dots + u v_k \]
        \item \textbf{Inductive Step:} Differentiating both sides once more with respect to $x$:
        \[ (uv)_{k+1} = \frac{d}{dx} [ (uv)_k ] \]
        Using the product rule on each term:
        \[ (uv)_{k+1} = \left( u_{k+1} v + u_k v_1 \right) + \binom{k}{1} \left( u_k v_1 + u_{k-1} v_2 \right) + \dots + \binom{k}{r} \left( u_{k+1-r} v_r + u_{k-r} v_{r+1} \right) + \dots \]
        Grouping like terms together:
        \[ (uv)_{k+1} = u_{k+1} v + \left[ 1 + \binom{k}{1} \right] u_k v_1 + \left[ \binom{k}{1} + \binom{k}{2} \right] u_{k-1} v_2 + \dots + \left[ \binom{k}{r} + \binom{k}{r-1} \right] u_{k+1-r} v_r + \dots + u v_{k+1} \]
        Using Pascal's Identity $\binom{k}{r} + \binom{k}{r-1} = \binom{k+1}{r}$:
        \[ (uv)_{k+1} = u_{k+1} v + \binom{k+1}{1} u_k v_1 + \binom{k+1}{2} u_{k-1} v_2 + \dots + u v_{k+1} \]
        This is the statement for $n = k+1$. Thus, by induction, the theorem is proved for all $n \in \mathbb{N}$.
    \end{itemize}

    \item[Step 3: Establish the second-order equation for $y = (\sin^{-1}x)^2$] \ \\
    Given:
    \[ y = (\sin^{-1}x)^2 \implies y(0) = 0 \]
    Differentiating once with respect to $x$:
    \[ y_1 = \frac{2\sin^{-1}x}{\sqrt{1-x^2}} \implies y_1(0) = 0 \]
    Squaring both sides:
    \[ y_1^2(1-x^2) = 4(\sin^{-1}x)^2 = 4y \]
    Differentiating again with respect to $x$:
    \[ 2y_1 y_2(1-x^2) - 2x y_1^2 = 4y_1 \]
    Dividing by $2y_1$ (since $y_1 \neq 0$):
    \[ y_2(1-x^2) - xy_1 = 2 \implies y_2(0) = 2 \]

    \item[Step 4: Differentiate $n$ times using Leibnitz's Theorem] \ \\
    Differentiating $y_2(1-x^2) - xy_1 - 2 = 0$ $n$ times:
    \[ \left[ y_{n+2}(1-x^2) + n y_{n+1}(-2x) + \frac{n(n-1)}{2} y_n (-2) \right] - \left[ y_{n+1}x + n y_n \right] = 0 \]
    \[ (1-x^2)y_{n+2} - 2nxy_{n+1} - n(n-1)y_n - x y_{n+1} - n y_n = 0 \]
    \[ (1-x^2)y_{n+2} - (2n+1)xy_{n+1} - n^2 y_n = 0 \]

    \item[Step 5: Find $(y_n)_0$] \ \\
    Setting $x=0$ in the above recurrence relation:
    \[ y_{n+2}(0) = n^2 y_n(0) \]
    Since $y_1(0) = 0$, all odd-order derivatives at $x=0$ are zero:
    \[ y_{2k+1}(0) = 0 \]
    For even-order derivatives, we use $y_2(0) = 2$:
    \[ y_4(0) = 2^2 \cdot y_2(0) = 2^2 \cdot 2 \]
    \[ y_6(0) = 4^2 \cdot y_4(0) = 4^2 \cdot 2^2 \cdot 2 \]
    \[ y_n(0) = 2 \cdot 2^2 \cdot 4^2 \cdot \dots \cdot (n-2)^2 \quad (\text{for even } n \ge 2) \]
\end{description}

Final Answer:
\[ 
\boxed{
y_n(0) = \begin{cases} 
2 \cdot 2^2 \cdot 4^2 \dots (n-2)^2, & \text{when } n \text{ is even} \\ 
0, & \text{when } n \text{ is odd} 
\end{cases}
} 
\]

%%%%%%%%%%%%%%%%%%%%%%%%%%%%%%%%%%%%%%%%%%%%%%%%%%  
% QUESTION_END  
%%%%%%%%%%%%%%%%%%%%%%%%%%%%%%%%%%%%%%%%%%%%%%%%%%

\end{document}